% concatenated from temporalcliques_rev.tex
\documentclass[11pt]{article}

% Set page size and margins
% Replace `letterpaper' with `a4paper' for UK/EU standard size
\usepackage[margin = 1in]{geometry}
\usepackage{setspace}


% Useful packages
\usepackage[T1]{fontenc}
\usepackage{charter} %font
\usepackage{euler} %mathfont
\usepackage{float}
\usepackage{soul} % has \ul for underlining
\usepackage{bbm} % has \mathbbm{1} for indicators
\usepackage{verbatim}
\usepackage{amsmath}
\usepackage{amsthm}
\usepackage{amssymb}
\usepackage{textcomp}
\usepackage{graphicx}
\usepackage[english]{babel}
\usepackage{tikz}
\usepackage{dblfloatfix}
\usepackage{csquotes}
\usepackage{titlesec}

\usepackage{cancel}


%enumeration modification from enumerate package%
\usepackage{enumerate}
\renewcommand\theenumi{(\roman{enumi})}
\renewcommand\theenumii{(\alph{enumii})}
\renewcommand{\labelenumi}{\rm (\roman{enumi})}
\renewcommand{\labelenumii}{\rm (\alph{enumii})}

%Bibliography stuff
% \usepackage[
% backend=biber,
% maxbibnames=99,
% maxcitenames=50,
% style=authoryear-comp,
% sorting=nyt
% ]{biblatex}
% \addbibresource{ref.bib}
% \usepackage{doi}


\allowdisplaybreaks %lets math break between pages

% writing algs
\usepackage{algorithm}
\usepackage{algpseudocode}

% formatting for theorems/lemmas
\newtheorem{theorem}{Theorem}[]
\newtheorem{conjecture}[theorem]{Conjecture}
\newtheorem{lemma}[theorem]{Lemma}
\newtheorem*{claim}{Claim}
\newtheorem{corollary}[theorem]{Corollary}
\newtheorem*{remark}{Remark}

%Other Formatting Nonsense

%Centers section titlees
\usepackage{sectsty}
\sectionfont{\centering}


% new commands for ease of use
\newcommand{\1}{\mathbf{1}} %Indicator
\newcommand{\Prob}{\mathbb{P}} %Probability
\newcommand{\prob}{\mathbb{P}} %Probability
\newcommand{\E}{\mathbb{E}} %Expected value
\newcommand{\ex}{\mathbb{E}} %Expected value
\newcommand{\N}{\mathbbm{N}} %Natural numbers
\newcommand{\R}{\mathbbm{R}} %Real numbers
\newcommand{\C}{\mathbbm{C}} %Complex numbers
\newcommand{\Q}{\mathbbm{Q}} %Rational numbers
\newcommand{\io}{\mbox{ i.o. }} % infinitely often
\newcommand{\as}{\mbox{ a.s. }} %almost surely
\newcommand{\id}{\ensuremath{\operatorname{Id}}} %identity
\newcommand{\lcm}{\ensuremath{\operatorname{lcm}}} %lcm
\newcommand{\bin}{\ensuremath{\operatorname{binomial}}} %binomial
\newcommand{\ber}{\ensuremath{\operatorname{bernoulli}}} %bernouilli
\newcommand{\unif}{\ensuremath{\operatorname{uniform}}} %uniform
\newcommand{\argmin}{\ensuremath{\operatorname{argmin}}} %argmin
\newcommand{\argmax}{\ensuremath{\operatorname{argmax}}} %argmax


\title{\bf{On the size of temporal cliques in subcritical random temporal graphs}
  \thanks{Luc Devroye acknowledges the support of NSERC grant A3450.
    G\'abor Lugosi acknowledges the support of Ayudas Fundación BBVA a
Proyectos de Investigación Científica 2021 and
the Spanish Ministry of Economy and Competitiveness grant PID2022-138268NB-I00, financed by MCIN/AEI/10.13039/501100011033,
FSE+MTM2015-67304-P, and FEDER, EU.).}
}


\author{
\bf{Caelan Atamanchuk} \\
Department of Mathematics and Statistics \\
McGill University \\
caelan.atamanchuk@gmail.com
\and
\bf{Luc Devroye} \\
School of Computer Science, McGill University, \\
  Montreal, Canada
lucdevroye@gmail.com
\and
\bf{G\'{a}bor Lugosi} \\
Department of Economics and Business, \\
Pompeu  Fabra University, Barcelona, Spain \\
ICREA, Pg. Lluís Companys 23, 08010 Barcelona, Spain \\
Barcelona Graduate School of Economics \\
gabor.lugosi@gmail.com
}


\begin{document}

\setstretch{1.2}

\maketitle


\begin{abstract}
  A \emph{random temporal graph} is an Erd\H{o}s-R\'enyi random graph $G(n,p)$, together with a random
  ordering of its edges. A path in the graph is called \emph{increasing} if the edges on the path appear
  in increasing order. A set $S$ of vertices forms a \emph{temporal clique} if for all $u,v \in S$, there is an increasing path from $u$ to $v$.
  Becker, Casteigts, Crescenzi, Kodric, Renken, Raskin and Zamaraev \cite{Becker2023} proved that if $p=c\log n/n$ for $c>1$, then,  with high probability, there
  is a temporal clique of size $n-o(n)$. On the other hand, for $c<1$,   with high probability,
  the largest temporal clique is of size $o(n)$.
  In this note, we improve the latter bound by showing that, for $c<1$, the largest temporal clique
  is of \emph{constant} size with high probability. 
\end{abstract}


\section{Introduction}\label{intro}

A \textit{temporal graph} $G = (V,E,\pi)$ is a graph $G=(V,E)$ together with an ordering $\pi:E \to \{1,\ldots,|E|\}$ on the edge set, interpreted as the times where the edges appear in the graph, often called the \textit{time stamps} of $G$. We say that an edge $e \in E$ \textit{precedes} an edge $e' \in E$ if $\pi(e) < \pi(e')$. A path from $u$ to $v$ is called $\textit{increasing}$ if each edge used in the path precedes the edge that is used after it, and we say that $v$ is \textit{reachable} from $u$ (or that $u$ can reach $v$) if an increasing path exists from $u$ to $v$. A set of vertices $S \subseteq V$ is called a \emph{temporal clique} if for all distinct vertices $u,v \in S$, there is an increasing path from $u$ to $v$ (and vice versa).

%One of the most interesting connectivity properties of temporal graphs is that, unlike in the case of ordinary graphs, reachability does not induce an equivalence relation on the vertices of the graph. Hence, a vertex can be in many distinct components at a single time.

In this note we discuss temporal graphs where $\pi$ is a uniform permutation on the edges and $G$ is an Erd\H{o}s-R\'{e}nyi random graph. The resulting temporal graph is called a \textit{random simple temporal graph}; \textit{RSTG} for short.
Motivated by modelling time-dependent propagation processes,
this model was introduced by Casteigts, Raskin, Renken, and Zamaraev \cite{Casteigts2023}.


One may generate RSTGs by a simple method: start with the complete graph $K_n$, then assign each edge an independent $\unif(0,1)$ random variable $(U_e : e \in E)$ and delete every edge with $U_e > p$. In this construction, we say that $e$ precedes $e'$ if $U_e < U_{e'}$. We also call the labels $U_e$ the time stamps. Importantly, creating i.i.d.\ uniform time stamps like this allows us to extend the notion of a temporal graph to infinite graphs which is needed for our analysis.

Casteigts, Raskin, Renken, and Zamaraev \cite{Casteigts2023} studied connectivity properties of RSTGs. They
identified the thresholds for different strengths of connectivity to
be in the region where $p = \frac{c\log(n)}{n}$ for some constant $c >
0$
  (Throughout the paper, $\log$ denotes natural logarithm).
Furthering this work, Broutin, Kamčev and Lugosi \cite{Broutin2023} identified the asymptotic lengths of the longest and shortest increasing paths in RSTGs with high probability for values of $p$ in this range. (We say that an event $E=E(n)$ happens with high probability if $\prob(E) \to 1$ as $n \to \infty$). 
Becker, Casteigts, Crescenzi, Kodric, Renken, Raskin and Zamaraev  \cite{Becker2023} identified $p = \frac{\log(n)}{n}$ as the threshold for the appearance of large temporal cliques. In particular, they showed that for every $\epsilon >0$,
when $p \ge \frac{(1+\epsilon)\log(n)}{n}$, then
there is a temporal clique of size $n-o(n)$ with high probability, while  if
$p \le \frac{(1-\epsilon)\log(n)}{n}$, then every temporal clique is of size $o(n)$ with high probability.

RSTGs are a natural way to model time-dependent processes on networks like social interactions and infection spread. A closely related model is the \textit{random gossip protocol} model, in which a sequence of edges $e_1,\ldots,e_k$ are chosen uniformly from the edges of $K_n$ and constructed a graph $G_{n,k}$. Increasing paths are defined as for temporal graphs. Papers studying this model include Moon \cite{Moon1972}, Boyd and Steele \cite{Boyd1979}
and Haigh \cite{Haigh1981}.
%Many more random graphs that are less related to the RSTG model have appeared in the discussion of gossip theory. One important one is the push-pull model, and the interested reader can look to Frieze \cite{Frieze1985}, Pittel \cite{Pittel1987} and Mocquard \cite{Mocquard2020} for an analysis.

For deterministic temporal graph models with random time stamps, see Chv\'{a}tal and Koml\'os \cite{Chvatal1971}, and Graham and Kleitman \cite{Graham1973}, 
Lavrov and Loh \cite{Lavrov2016} and Angel, Ferber, Sudakov and Tassion \cite{Angel2020}.

Our contribution is summarized in the following theorem.
It shows that in the subcritical regime $p= c\log n/n$ with $c<1$,
the size of the largest temporal clique is not only $o(n)$ but in fact, of size $O(1)$,
improving the upper bound of \cite{Becker2023}. This reveals a quite spectacular
phase transition around $p= \log n/n$, since for $c>1$, there is a temporal clique
of size $n-o(n)$. The behavior of the size of the largest temporal clique near the critical
regime remains an intriguing research problem.

\begin{theorem}\label{main}
Let $p = \frac{c\log(n)}{n}$, and let $G$ be an RSTG with edge probability $p$. If $c \in (0,1)$, then the largest temporal clique in $G$ is of size at most $\lceil \frac{1}{1-c}+1 \rceil$ with high probability.
\end{theorem}

Note that for $c \leq \frac{1}{2}$, Theorem \ref{main} asserts that $G$ has
no temporal clique of size $4$. This bound can't be improved, since 
for $p = \omega(\frac{1}{n})$ the static Erd\H{o}s-R\'{e}nyi graph contains a triangle with high probability. Moreover, every triangle is trivially a temporal clique of size $3$.
We conjecture that the upper bound of Theorem \ref{main} is sharp for all $c\in (0,1)$.

The proof of Theorem \ref{main} is based on relating the number of vertices that are reachable by monotone paths from a typical vertex to the total progeny of a certain ``temporal'' branching process.
We utilize the temporal branching process bounds to assert that the number of vertices that a collection of $m \geq \lceil 1 + \frac{1}{1-c}\rceil$ vertices can reach is small enough so that the chance of them forming a component unlikely enough that expected number of components of size $m$ tends to zero.

\section{Temporal branching processes}\label{temptree}

We begin this section by introducing temporal branching processes that are the key tool
in the proof of Theorem \ref{main}.
% (see \cite{Athreya1972} for the definition of ordinary branching processes).
We only focus on branching processes with binomial offspring distribution as this is the degree distribution of a typical vertex in an Erd\H{o}s-R\'enyi random graph.
One way to generate these processes is as follows.
Start with an infinite rooted $n$-ary tree, add an independent $\unif(0,1)$ time stamp $U_e$
to every edge. Delete any edge with $U_e > p$. This decomposes the tree into a forest
and we only focus on the component that contains the root vertex.
We say that a vertex $v$ is reachable from the root if the unique path from the root to $v$ in the infinite $n$-ary tree has edge labels $U_e \leq p$ for each edge on the path and these labels are increasing on the path.
The subtree consisting of only vertices reachable from the root is a temporal branching process with a $\bin(n,p)$ offspring distribution. Throughout the rest of the paper, $T$ is always an infinite $n$-ary tree with such a labelling on the edges. The following sequence of results provides us with the necessary upper bounds for the size of the reachable set of $T$.


\begin{lemma}\label{treesep}
Let $P_1,\ldots,P_q$ be a finite collection of distinct infinite paths in $T$, and let $(X_k)_{k \geq 0}$ be a random walk down the tree, that is, $X_0$ is the root and $X_k$ is uniformly distributed over the children of $X_{k-1}$ for all $k \geq 1$. Then, $\prob(\tau \geq \ell) \leq \frac{q}{n^\ell}$, where
$$\tau = \max\{ k > 0 : X_0,\ldots,X_k \text{ coincides with one of the $q$ paths}\}.$$
\end{lemma}


\begin{proof}
If $\tau \geq \ell$, then $X_0,\ldots,X_\ell$ coincides with one of the $P_1,\ldots,P_q$. Since $X_k$ is uniform over the children of $X_{k-1}$ and the tree is $n$-ary, the result follows from the union bound yields.
\end{proof}


\begin{lemma}\label{treesurvival}
Let $T^*$ be the set of all vertices that are reachable from the root of $T$. If $v_1,\ldots,v_q$ are uniform vertices chosen from the $\ell$-th generation of $T$, then
$$\prob(v_1,\ldots,v_q \in T^*) \leq \frac{(q-1)!e^{npq}}{n^{\ell q}}.$$
Furthermore, when $\ell \geq (np)^4$ and $np \to \infty$ as $n \to \infty$,
$$\prob(v_1,\ldots,v_q \in T^*) = O\left(\frac{(q-1)!C^{q-1}}{n^{\ell q}}\right),$$
for some $C > 0$.
\end{lemma}


\begin{proof}
Suppose that $v_1,\ldots,v_r \in T^*$ and let $T'$ be a subtree of $T$ consisting of $r$ distinct infinite paths starting at the root through $v_1,\ldots,v_r$. Let $(X_k)_{k \geq 0}$ be a random walk down the tree (independent of $v_1,\ldots,v_r$) and let $\tau$ be as in Lemma \ref{treesep}. If $\tau = j$, then there are $\ell-j$ edges left that need to both exist and be increasing to have $X_\ell \in T^*$. Hence, if $V_r = \{v_1 , \ldots , v_r \in T^*\}$,
$$\prob(X_\ell \in T^* | V_r \cap \{\tau = j\}) = p^{\ell-j}\prob\big(\big\{\text{a path of length $\ell$ is increasing}\big\} \big|\big\{ \text{first $j$ edges are increasing}\big\}\big) = \frac{p^{\ell-j}j!}{\ell!}.$$
Combining this with Lemma \ref{treesep} yields
\begin{align}\label{sum}
\prob(X_\ell \in T^* | V_r)  &
\leq \sum_{j=0}^{\ell-1} \prob(\tau = j)\prob(X_\ell \in T^* | V_r \cap \{\tau = j\}) 
+ \prob(\tau \ge \ell)\prob(X_\ell \in T^* | V_r \cap \{\tau \ge \ell\}) 
 \nonumber \\
& \leq \sum_{j=0}^\ell \left(\frac{r}{n^j}\right)\frac{p^{\ell-j}j!}{\ell!} = \frac{r}{n^\ell}\sum_{j=0}^\ell\frac{(np)^{\ell-j}j!}{\ell!}.
\end{align}
To get the first bound we use the fact that $\frac{j!}{\ell!} \leq \frac{1}{(\ell-j)!}$ to get 
$$\prob(X_\ell \in T^*|V_r) \leq \frac{re^{np}}{n^\ell}.$$
Then, since $X_\ell$ is distributed uniformly across the $\ell$-th generation,  applying the above inequality repeatedly,
\begin{align*}
\prob(v_1,\ldots,v_q \in T^*) &= \prob(v_1 \in T^*) \cdot \prob(v_2 \in T^* | V_1) \cdots \prob(v_q \in T^* | V_{q-1}) \\
&\leq \frac{p^\ell}{\ell!}\prod_{r=1}^{q-1}\left(\frac{re^{np}}{n^\ell}\right) = \left(\frac{(np)^{\ell}}{\ell!}\right)\frac{(q-1)!e^{np(q-1)}}{n^{\ell q}} \leq \frac{(q-1)!e^{npq}}{n^{\ell q}}.
\end{align*} 
For the second inequality, we split the sum in (\ref{sum}) in two separate pieces
\begin{align*}
    A = \sum_{j=0}^{\ell - \sqrt{\ell}}\frac{(np)^{\ell - j}j!}{\ell!}, \quad \text{and} \quad B = \sum_{j=\ell -\sqrt{\ell}+1}^\ell \frac{(np)^{\ell-j}j!}{\ell!} = \sum_{k=0}^{\sqrt{\ell}-1} \frac{(np)^k}{\ell \cdot (\ell-1) \cdots (\ell-k+1)}.
\end{align*}
The first term may be bounded by
\begin{align*}
A &\leq \sum_{k=\sqrt{\ell}}^\ell \frac{(np)^k}{k!} \leq \left|e^{np}- \sum_{k=0}^{\sqrt{\ell}-1} \frac{(np)^k}{k!}\right| \leq \frac{e^{np}(np)^{\sqrt{\ell}}}{(\sqrt{\ell})!}~,
\end{align*}
where the second inequality follows from the Lagrange form of the remainder in Taylor's theorem.

For the second term, since $\ell \cdot (\ell-1) \cdots (\ell - k + 1) \geq \ell^{k}(1-\frac{1}{\sqrt{\ell}})^{\sqrt{\ell}}$ for all $0 \leq k \leq \sqrt{\ell}$, we have that
\begin{align*}
    B &\leq \frac{1}{\left(1-\frac{1}{\sqrt{\ell}}\right)^{\sqrt{\ell}}}\sum_{k=0}^{\sqrt{\ell}-1} \left(\frac{np}{\ell}\right)^k \leq \frac{1}{\left(1-\frac{np}{\ell}\right)\left(1-\frac{1}{\sqrt{\ell}}\right)^{\sqrt{\ell}}} = O(1),
\end{align*}
when $\ell \geq (np)^4$. Combining both the bounds along with Stirling's approximation, we conclude that there is some $C > 0$ such that
$$
\prob(X_\ell \in T^* | V_r) 
\leq \frac{r}{n^\ell}\left(\frac{e^{np}(np)^{\sqrt{\ell}}}{\sqrt{\ell}!} + O(1)\right) \leq \frac{r}{n^\ell}\left(\frac{e^{np}(np)^{(np)^2}e^{(np)^2}}{(np)^{2(np)^2}}+O(1)\right) \leq \frac{Cr}{n^\ell},$$
when $\ell \geq (np)^4$ and $np \to \infty$. Proceeding exactly as we did for the first inequality,
\begin{align*}
\prob(v_1,\ldots,v_q \in T^*) &= \prob(v_1 \in T^*) \cdot \prob(v_2 \in T^* | V_1) \cdots \prob(v_q \in T^* | V_{q-1}) \\
&\leq \frac{p^\ell}{\ell!}\frac{(q-1)!C^{q-1}}{n^{\ell (q-1)}} \leq \frac{(np)^{(np)^4}}{(np)^4!}\frac{(q-1)!C^{q-1}}{n^{\ell q}} = O\left( \frac{(q-1)!C^{q-1}}{n^{\ell q}}\right),
\end{align*} 
where in the final bound we use the fact that $\frac{x^{x^4}}{(x^4)!} = o(1)$ as $x\to \infty$, which is an immediate consequence of Stirling's approximation.
\end{proof}

We may use Lemma \ref{treesurvival} to bound the moments of the number of vertices reachable in a particular generation. For $\ell \geq 0$, denote by $Z_\ell$ the number of vertices in $T$ reachable from the root in the $\ell$-th generation. 

\begin{corollary}\label{reversejensen}
For all integers $\ell \ge 0$ and $q \ge 1$,
$\ex[Z_\ell^q] \leq (q -1)!e^{npq}$. Furthermore, when $\ell \geq (np)^4$ and $np \to \infty$, there is a constant $C > 0$ such that $\ex[Z_\ell^q] = O((q-1)!C^{q-1})$.
\end{corollary}


\begin{proof}
Denoting by $S_\ell$ the set of $n^{\ell}$ vertices in the $\ell$-th generation of $T$, we may write
$Z_\ell = \sum_{v\in S_{\ell}} \mathbf{1}_{\{v\in T^*\}}$. Then
\[
  \ex [Z_\ell^q] = n^{\ell q}\prob( v_1,\ldots,v_q \in T^*)~,
\]
where $v_1,\ldots,v_q$ are independent vertices chosen uniformly at random from $S_\ell$.
Combining this with Lemma \ref{treesurvival} implies the stated bounds.
\end{proof}



The next bound will control the number of vertices that a typical vertex in a simple random temporal graph can reach, further allowing us to control the size of temporal cliques.

\begin{theorem}\label{sizebound} Let $T^*$ be the set of vertices in $T$ that are reachable from the root and suppose that $np \to \infty$. Then, for any integer $q \geq 1$, there is a constant $c(q)$ such that
$$\ex[|T^*|^q] \leq c(q)(np)^{4q}e^{npq}.$$
\end{theorem}


\begin{proof}
%Using H\"{o}lder's inequality along with the identity $\ex[X^q] = \int qt^{q-1}\prob(X > t) dt$ yields
Observe that
\begin{align*}
    \ex|T^*|^q & = \ex\left(\sum_{i=0}^{(np)^4}Z_i + \sum_{i > (np)^4} Z_i\right)^q  \\
& \le 2^{q-1} \ex\left(\sum_{i=0}^{(np)^4}Z_i \right)^q + 2^{q-1} \ex\left(\sum_{i > (np)^4} Z_i\right)^q
\quad \text{(by Jensen's inequality)} \\
&
\leq \underbrace{2^{q-1}((np)^4+1)^{q-1}\ex\left[ \sum_{i=0}^{(np)^4} Z_i^q \right]}_{:= I} + \underbrace{2^{q-1}q\sum_{t > 0}t^{q-1}\prob\left(\sum_{i > (np)^4} Z_i \geq t\right)}_{:= II}~,
\end{align*}
where in the last step we used Jensen's inequality to bound the first term and 
the identity $\ex[X^q] = \int qt^{q-1}\prob(X > t) dt$ to bound the second.
The first inequality of Corollary \ref{reversejensen} may be used to bound 
the expectation in term $I$,
 as 
$$\ex\left[ \sum_{i=0}^{(np)^4} Z_i^q\right] \leq 
((np)^{4}+1)
\left(\sup_{i \geq 0}\ex[Z_i^q]\right) \leq
((np)^{4}+1)
(q - 1)!e^{npq}.$$
To bound $II$, we may write
\begin{align*}
t^{q-1}\prob\left( \sum_{i > (np)^4} Z_i \geq t\right) &\leq t^{q-1}\prob\left( Z_{(np)^4 + \log t} > 0\right) + t^{q-1}\prob\left(\bigcup_{i=(np)^4 + 1}^{(np)^4 + \log t} \left\{Z_i \geq \frac{t}{\log t}\right\}\right) \\
&\leq \underbrace{t^{q-1}\prob\left( Z_{(np)^4 + \log t} > 0\right)}_{:=III} + \underbrace{t^{q-1}(\log t)\max_{(np)^4 < i \leq (np)^4 + \log t} \prob\left(Z_i \geq \frac{t}{\log t}\right)}_{:= IV},
\end{align*}
and we can bound the two terms separately. To bound $III$, note that at level $(np)^4+\log t$, there are $n^{(np)^4 + \log t}$ vertices, and they are each reachable with probability $p^{(np)^4 + \log t}/((np)^4+ \log t)!$. Thus, by Stirling's approximation,
\begin{align*}
III &\leq \frac{t^{q-1}(np)^{(np)^4 + \log t}}{((np)^4 + \log t)!} \\   
    &\leq \frac{t^{q-1}(enp)^{(np)^4 + \log t}}{(np)^{4(np)^4 + 4\log t}} \\
    &  = t^{q-1}\left( \frac{e}{(np)^3} \right)^{(np)^4 + \log t} \\
    &= t^{q-3\log(np)}\left(\frac{e}{(np)^3}\right)^{(np)^4}
\end{align*}
for any $t \geq 0$. In particular, this implies that $III$ is summable and converges to 0 when $np \to \infty$. Applying Markov's inequality and the second inequality in Corollary \ref{reversejensen} gives
$$IV \leq O\left( \frac{\log^{k+1} (t) ( k -1 )!
C^{k-1}
}{t^{k-q+1}}\right),$$
for any positive integer $k$ and $t \geq 0$. Choosing $k = q+1$ results in $IV$ being summable and bounded above by a constant depending only on $q$. Grouping up all that only depends on $q$ and upper bounding by some dominating constant $c(q)$ we get
$$\ex|T^*|^q \leq I + III + IV \leq c(q)(np)^{4q}e^{npq}.$$
\end{proof}



\section{Proof of Theorem \ref{main}}

We are now prepared to prove Theorem \ref{main}.
For labelled vertices $\{1,\ldots,m\}$ to form a temporal clique in an RSTG they need to all be reachable from one another. This can happen if and only if for all distinct $u,v \in \{1,\ldots,m\}$, there is a vertex $w$ that can reach $v$ with only edges that have time stamps above $p/2$, and is reachable from $u$ with only edges that have time stamps below $p/2$. With this in mind, for any $0\le a < b\le p$, we define $G_{[a,b]}$ to be the subgraph obtained from $G$ by only keeping edges with time stamps in $[a,b]$. Set $A_1,\ldots,A_m$ to be the collection of all vertices that are reachable from $1,\ldots,m$ in $G_{[0,p/2]}$ and $B_1,\ldots,B_m$ to be the collection of all vertices in $G_{[p/2,p]}$ that can reach $1,\ldots,m$. With this new notation, we can say that $\{1,\ldots,m\}$ form a temporal clique if for all $i,j \in \{1,\ldots,m\}$ distinct, the set $A_i \cap B_j$ is nonempty. 

It is important to note that $G_{[p/2,p]}$ and $G_{[0,p/2]}$ are identically distributed RSTGs, but are not independent. Observe that for any $0\le a < b \le p$, the RSTG $G_{[a,b]}$ is determined by the binary vector $X = (X_1,\ldots,X_{\binom{n}{2}})$
defined by $X_i = \1_{e_i \in G_{[a,b]}}$ (for some enumeration of the edges of $K_n$) and a random permutation $O_{[a,b]}$ of
$[\binom{n}{2}]$ that denotes the relative orderings of the edge labels. In the next lemma we consider certain functionals
of $G_{[a,b]}$, represented by $X$ and $O_{[a,b]}$. More precisely, such a functional is of the form
$f:\{0,1\}^{\binom{n}{2}}\times \mathrm{Sym}(\binom{n}{2}) \to \R$, where $\mathrm{Sym}(\binom{n}{2})$ is the set of permutations
of $[\binom{n}{2}]$.
The next lemma deals with the dependence between two subgraphs $G_{[a,b]}$ and $G_{[c,d]}$, for some $0 \leq a < b < c < d \leq p$.


\setcounter{theorem}{5}

\begin{lemma}\label{dependence}
Let $G$ be an RSTG with vertices labelled $\{1,\ldots,n\}$, and let $0 \leq a < b < c < d \leq p$. Set $X_i = \1_{e_i \in G_{[a,b]}}$, $Y_i = \1_{e_i \in G_{[c,d]}}$ for some enumeration of the edges of $K_n$, $e_1,\ldots,e_{\binom{n}{2}}$, and let $X = (X_1,\ldots,X_{\binom{n}{2}})$ and $Y = (Y_1,\ldots,Y_{\binom{n}{2}})$. Let $O_{[a,b]}$ and $O_{[c,d]}$ be the permutations that denote the relative orderings of the edges in the two graphs. 
Let
$f,g:\{0,1\}^{\binom{n}{2}}\times \mathrm{Sym}(\binom{n}{2}) \to \R$ be such that
 $f(x_1,\ldots,x_{\binom{n}{2}},s)$ and $g(x_1,\ldots,x_{\binom{n}{2}},s)$ are two non-decreasing functions in $x_1,\ldots,x_{\binom{n}{2}}$ for any fixed $s\in  \mathrm{Sym}(\binom{n}{2})$. 
Then
$$\ex\left[f\left(X,O_{[a,b]}\right)g\left(Y,O_{[c,d]}\right)\right] \leq \ex\left[f\left(X,O_{[a,b]}\right)\right] \ex\left[g\left(Y,O_{[c,d]}\right)\right].$$
In particular,
$$\ex\left[ |A_1|^q \cdots |A_m|^q \cdot |B_1|^q \cdots |B_m|^q \right] \leq \ex\left[|A_1|^q \cdots |A_m|^q \right]\ex\left[
|B_1|^q \cdots |B_m|^q\right] = \ex\left[|A_1|^q \cdots |A_m|^q\right]^2,$$
for all $q \geq 0$ and $A_1,B_1,\ldots,A_m,B_m$ defined as above.
\end{lemma}

\begin{proof}
Conditioned on $Z= (X_1,Y_1,\ldots,X_{\binom{n}{2}},Y_{\binom{n}{2}})$, all of the randomness of $f\left(X,O_{[a,b]}\right)$ and $g\left(Y,O_{[c,d]}\right)$ comes from the random relative orderings. Since $a < b < c < d$, the two random variables $O_{[a,b]}$ and $O_{[c,d]}$ are independent, which implies that $f\left(X,O_{[a,b]}\right)$ and $g\left(Y,O_{[c,d]}\right)$ must also be conditionally independent on $Z$. Hence, by the tower property of conditional expectation,
\begin{align*}
\ex\left[f\left(X,O_{[a,b]}\right)g\left(Y,O_{[c,d]}\right)\right] &= \ex\left[\ex\left[f\left(X,O_{[a,b]}\right)g\left(Y,O_{[c,d]}\right) \Big| Z \right]\right] \\
&= \ex\left[\ex\left[f\left(X,O_{[a,b]}\right) \Big| Z \right]\ex\left[g\left(Y,O_{[c,d]}\right) \Big| Z \right]\right]. \\
&= \ex\left[\ex\left[f\left(X,O_{[a,b]}\right) \Big| X \right]\ex\left[g\left(Y,O_{[c,d]}\right) \Big| Y\right]\right],
\end{align*}
where the final equality just follows from the fact that, once we condition on $X$, knowing $Y$ tells us nothing about $f(X,O_{[a,b]})$ and vice versa. The random variables
$$\ex\left[f\left(X,O_{[a,b]}\right) \Big| X = (z_1,\ldots,z_{\binom{n}{2}}) \right], \ \text{and} \ \ex\left[g\left(Y,O_{[c,d]}\right) \Big| Y = (z_1,\ldots,z_{\binom{n}{2}}) \right]$$
are non-decreasing functions in $z_1,\ldots,z_{\binom{n}{2}}$ by the definitions of $f$ and $g$. Furthermore, the collection of random variables $\{X_i : i \in \{1,\ldots,\binom{n}{2}\}\} \cup \{Y_i : i \in \{1,\ldots,\binom{n}{2}\}\}$ are negatively associated (this can be seen by combining Proposition 7 and Lemma 8 from Dubhashi and Ranjan \cite{Dubhashi1998}). With this, applying the tower property again gives
\begin{align*}
\ex\left[f\left(X,O_{[a,b]}\right)g\left(Y,O_{[c,d]}\right)\right]  \leq \ex\left[f\left(X,O_{[a,b]}\right)\right]\ex\left[g\left(Y,O_{[c,d]}\right)\right].   
\end{align*}
Observing that $|A_1| \cdots |A_m|$ and $|B_1| \cdots |B_m|$ satisfy the conditions of the first statement and are identically distributed is enough to complete the proof of the second inequality.
\end{proof}


The next lemma acts as a bridge between RSTGs and the temporal branching processes explored in the previous section. The idea behind the proof hinges on the fact that the sizes of $\bin(n,p)$ branching processes upper bounds the sizes of neighbourhoods around vertices in an Erd\H{o}s-R\'{e}nyi graph, though formalizing this idea takes some work. Equipped with this and Theorem \ref{sizebound}, the proof of Theorem \ref{main} is reduced to a routine use of the first-moment method.

\begin{lemma}\label{stoch}
$|A_1|$ is stochastically dominated by $|T^*|$, where $T^*$ is the set of vertices reachable from the root in a temporal branching process $T$ with offspring distribution $\bin(n,p/2)$. In particular, $\ex[|A_1|^q] \leq \ex[|T^*|^q]$ for all $q \geq 0$.
\end{lemma}

\begin{proof}
We can determine $A_1$ via the foremost tree algorithm from
Casteigts, Raskin, Renken, and Zamaraev \cite{Casteigts2023}.
  The algorithm builds a tree recursively, building a tree of
  increasing
  paths starting from an arbitrary
  vertex. The algorithm is defined  as follows:

\begin{itemize}
\item
    Initialize with $\tau_0 = 1$ and $G_0$ as the single vertex
    labelled $1$.
 \item
  While $\tau_k \leq p/2$, set $\tau_{k+1}$ to be the smallest time
stamp of an edge   connecting vertices in $G_k$ with vertices outside of $G_k$
that is
larger than $\tau_k$.
\item
  If $\tau_{k+1} \leq p/2$, add the corresponding edge $e_{k+1}$ and
vertex $v_{k+1}$ to obtain $G_{k+1}$.
\item
    If $\tau_k>p/2$, the algorithm terminates and outputs
$G_{k}$.
\end{itemize}

Note that $|A_1|$
  equals the number of vertices of the resulting tree $G_k$,
that is,
\begin{equation}\label{time}
|A_1| = \inf\left\{k \geq 0 : \tau_k > \frac{p}{2}\right\}~.
\end{equation}
This foremost tree algorithm can also be run on the tree $T$ as a way
to generate $T^*$ with the same procedure, and we denote the sequence
of timestamps in this graph as $\tau^*_k$. Additionally, we denote by
$E_k$ and $E^*_k$ the collection of all viable edges that could be
added during step $k$, that is, all edges from $G_k$ to $K_n$ that, if added, keep the graph $G_{k+1}$ as an increasing tree (all vertices reachable from the root). Note that by the definition of the algorithm, every edge in $E_k$ must have a time stamp that is at least $\tau_k$ and similarly for $E_k^*$ and $\tau^*_k$. Moreover, the time stamps of edges in $E_k$ and $E^*_k$ are uniformly distributed on $[\tau_{k-1},1]$ and $[\tau_{k-1}^*,1]$ respectively. By means of a direct inductive coupling we show that $\tau$ stochastically dominates $\tau^*$, $|E_k^*|$ stochastically dominates $|E_k|$, and hence $|T^*|$ must stochastically dominate $|A_1|$ by the characterization of (\ref{time}).
%(see, eg., Thorisson \cite{Thorisson1995} or Thorisson \cite{Thorisson2000} for the relevant definitions as well as information on the use of couplings for stochastic domination results). In the proof below, $(U_{ij} : i,j \geq 1)$ is taken to be a countable collection of random variables uniform on $[0,1]$.

The base case of the induction is easy to see. By definition $|E_1| = n-1$, $|E_1^*| = n$,
  $\tau_1 \sim \min_{1 \leq i \leq n-1} U_{1,i}$, and $\tau_1^* \sim \min_{1 \leq n} U_{1,i}$. Thus, just using the same uniforms to generate both $\tau_1$ and $\tau_1^*$ is enough. Now suppose that there is some probability space $(\Omega_{k-1},\mathcal{F}_{k-1},\prob_{k-1})$ and random variables distributed as $|E_{k-1}|,|E_{k-1}^*|,\tau_{k-1},\tau^*_{k-1}$ (we just use the same symbols to denote these random variables) such that $|E_{k-1}|(\omega) \leq |E_{k-1}^*|(\omega)$ and $\tau^*_{k-1}(\omega) \leq \tau_{k-1}(\omega)$ for all $\omega \in \Omega_{k-1}$. In the $(k-1)$-th step of the algorithm we added a new vertex to both graphs, resulting in $n-k$ possible new edges to $G$ and $n$ edges to $T$ for the $k$-th step. Hence, since we cannot add edges that are below $\tau_{k-1}$ and $\tau_{k-1}^*$ respectively
$$|E_k| \sim |E_{k-1}| + \bin(n-k,1-\tau_{k-1}), \quad |E_k^*| \sim |E_{k-1}^*| + \bin(n,1-\tau_k^*),$$
and, recalling the distribution of time stamps in $E_k$ and $E_k^*$,
$$\tau_k \sim \tau_{k-1} + (1-\tau_{k-1})\min_{1 \leq i \leq |E_k|} U_{k,i}, \quad \tau_k^* \sim \tau^*_{k-1} + (1-\tau^*_{k-1})\min_{1 \leq i \leq |E_k^*|} U_{k,i}.$$
Let $(\Omega_k, \mathcal{F}_k,\prob_k)$ be the product of $(\Omega_{k-1},\mathcal{F}_{k-1},\prob_{k-1})$ with $(\Omega', \mathcal{F}',\prob')$, the probability space of $\binom{n}{2}+n$ independent uniform random variables, $(U_{k,i} : 1 \leq i \leq \binom{n}{2}+n)$. Here we can couple the binomial random variables by generating them as $\sum_{i=1}^k \1_{\{U_{k,i} \leq (1-\tau_{k-1})\}}$ and $\sum_{i=1}^k \1_{\{U_{k,i} \leq (1-\tau^*_{k-1}\}}$ respectively. Then, if we generate $|E_k|$ and $|E_k^*|$ with these binomials, by the inductive hypothesis, it must be the case that $|E_k|(\omega) \leq |E_k^*|(\omega)$ for all $\omega \in \Omega_k$. Similarly, using the uniforms $(U_{k,i} : n+1 \leq i \leq |E_k^*|)$ to generate both $\tau_k$ and $\tau_k^*$ according to their distributions results in also having $\tau_k^*(\omega) \leq \tau_k(\omega)$ for all $\omega \in \Omega_k$.
\end{proof}

With the lemmas out of the way we now prove our main result.

\begin{proof}[Proof of Theorem \ref{main}]
Let $m \geq 0$ and let $A_1,\ldots,A_m,B_1,\ldots,B_m$ be as defined in the beginning of this section. Let $N$ be the number of temporal cliques of size $m$ in $G$. Then, if we take $(v_{ij})_{i,j = 1}^m$ to be independently and uniformly chosen random vertices from the labelled set $\{1,\ldots,n\}$, we may apply Lemma $\ref{dependence}$ to get that
\begin{align*}
\ex[N] &\leq n^m\prob(\{1,\ldots,m\} \text{ form a temporal clique}) \\
&= n^{m+m(m-1)}\prob(v_{ij} \in A_i, v_{ij} \in B_j \ \forall i \ne j) \\
&= n^{m+m(m-1)}\ex\left[ \prob\Big(v_{ij} \in A_i, v_{ij} \in B_j \ \forall i \ne j \Big| |A_1|,\ldots,|A_m|,|B_1|,\ldots,|B_m| \Big)  \right] \\
&\leq n^{m+m(m-1)}\ex\left[ \left( \frac{|A_1|}{n} \right)^{m-1} \cdots \left( \frac{|A_m|}{n} \right)^{m-1} \left( \frac{|B_1|}{n} \right)^{m-1} \cdots \left( \frac{|B_m|}{n} \right)^{m-1} \right] \\
&\leq n^{m+m(m-1)}\ex\left[ \left( \frac{|A_1|}{n} \right)^{m-1} \cdots \left( \frac{|A_m|}{n} \right)^{m-1}\right]^2 \\
&\leq n^{m-m(m-1)}\ex[|A_1|^{m-1} \cdots |A_m|^{m-1}]^2.
\end{align*}
Applying H\"{o}lder's inequality along with Lemma \ref{stoch} and
Theorem \ref{sizebound}
  applied for the probability $p/2 = c log(n)/(2n)$,
gives us the upper bound 
\begin{align*}
  \ex[N]
  &
    \leq n^{m-m(m-1)}\ex[|A_1|^{m(m-1)}]^2 \leq
         \kappa_m \left(\frac{c}{2}\log(n)\right)^{8m(m-1)}n^{m-m(m-1)+c(m(m-1))}
\end{align*}
  where $\kappa_m$ is a constant depending on $m$ only.
If $m-m(m-1)+cm(m-1) < 0$, then $\ex[N] \to 0$ as $n \to \infty$. Rearranging, this inequality is equivalent to $m \geq \lceil \frac{1}{1-c}+1\rceil$ as $m$ is an integer.
\end{proof}



\bibliographystyle{plain}
\bibliography{ref}
	


\end{document}
